\begin{solucion}
    \textbf{Demostración 1.}

    Definamos $\left(\binom{[n]}{k}\right)$, como el conjunto de todos los posibles k-combinaciones del multiconjunto:
    \begin{align*}
        \{
            1\cdot\infty, 2\cdot\infty, \ldots, n\cdot\infty
        \}
    \end{align*}
    Es decir:
    \begin{align*}
        \left(\binom{[n]}{k}\right) = \left\{
            \{k_1\cdot 1, k_2\cdot 2, \ldots, k_n\cdot n\}:k_i \in \mathbb N \text{ para todo $i$ } \text{ y } \sum_{i=1}^n k_i = k
        \right\}
    \end{align*}
    Note que $k_i$ puede ser cero para algunos $i = 1, \ldots, n$. De esta forma:
    \begin{align*}
        \left|\left(\binom{[n+2]}{m}\right)\right|= \left|
            \{
                \{k_1\cdot 1, k_2\cdot 2, \ldots, k_{n+2}\cdot (n+2)\}:k_i \in \mathbb N \text{ para todo $i$ } \text{ y } \sum_{i=1}^{n+2} k_i = m
            \}
        \right|
    \end{align*}
    \begin{align*}
        \left|
           \bigcup_{k=0}^m \left(\binom{[n+1]}{k}\right)
        \right| = \sum_{k=0}^m \left|\left(\binom{[n+1]}{k}\right)\right|
    \end{align*}
    Lo cual es posible pues para $i\neq j$, $\left(\binom{[n+1]}{i}\right)\cap\left(\binom{[n+1]}{j}\right)=\emptyset$, pues $\left(\binom{[n+1]}{i}\right)$ contiene multiconjuntos de tamaño $i$ y $\left(\binom{[n+1]}{j}\right)$ contiene multiconjuntos de tamaño $j$.

Asi definimos la función $f$ como:
\begin{align*}
    f:\left(\binom{[n+2]}{m}\right) &\to \bigcup_{k=0}^m \left(\binom{[n+1]}{k}\right)\\
    \{k_i\cdot i\}_{i=1}^{n+2} &\mapsto \{k_i\cdot i\}_{i=1}^{n+1} 
\end{align*}

Note que si $k_i = 0$ entonces $i \notin \left(\binom{[n+1]}{k}\right)$. De este modo si $k_i = 0$ para todo $i \leq n+1$ entonces $k= k_1+k_2+\ldots+k_{n+1} = 0$ y por tanto:
\begin{align*}
    \{0\cdot 1, 0\cdot 2, \ldots, 0\cdot n+1\} = \emptyset
\end{align*}
Por lo que:
\begin{align*}
    \left(\binom{[n+1]}{0}\right) = \left\{
        \{0\cdot 1, 0\cdot 2, \ldots, 0\cdot n+1\}
    \right\}=\left\{
        \emptyset
    \right\}
\end{align*}        
$f$ es también definida, por lo que veamos que $f$ es inyectiva:
\begin{align*}
    f(\{k_i\cdot i\}_{i=1}^{n+2}) = f(\{k'_i\cdot i\}_{i=1}^{n+2})\quad \implies \quad  \{k_i\cdot i\}_{i=1}^{n+1} = \{k'_i\cdot i\}_{i=1}^{n+1}
\end{align*}
Luego $k_i = k'_i$ para todo $i=1,\ldots,n+1$. Por lo que:
\begin{align*}
    \{k_i\cdot i\}_{i=1}^{n+2} = \{k'_i\cdot i\}_{i=1}^{n+2}
\end{align*}
Por lo que $f$ es inyectiva.

Ahora veamos que $f$ es sobreyectiva. Sea 
\begin{align*}
    \{k_i\cdot i\}_{i=1}^{n+1} \in \bigcup_{k=0}^m \left(\binom{[n+1]}{k}\right)
\end{align*}
Entonces existe un $k$ tal que:
\begin{align*}
    \{k_i\cdot i\}_{i=1}^{n+1} \in \left(\binom{[n+1]}{k}\right)
\end{align*}
con $k_1 + k_2 + \ldots + k_{n+1} = k$. Definamos $k_{n+2} = m - k$ entonces:
\begin{align*}
    \sum_{i=1}^{n+2} k_i = k_1 + k_2 + \ldots + k_{n+1} + k_{n+2} = k + (m - k) = m
\end{align*}
por lo que:
\begin{align*}
    \{k_i\cdot i\}_{i=1}^{n+2} = \{k_1\cdot 1, k_2\cdot 2, \ldots, k_{n+1}\cdot (n+1), k_{n+2}\cdot (n+2)\} \in \left(\binom{[n+2]}{m}\right)
\end{align*}
De este modo para $\{k_i\cdot i\}_{i=1}^{n+1} \in \bigcup_{k=0}^m \left(\binom{[n+1]}{k}\right)$ existe $\{
k_i\cdot i\}_{i=1}^{n+2} \in \left(\binom{[n+2]}{m}\right)$, con $k_{n+2} = m - k$, tal que:
\begin{align*}
    f(\{k_i\cdot i\}_{i=1}^{n+2}) = \{k_i\cdot i\}_{i=1}^{n+1}
\end{align*}
Por lo que $f$ es sobreyectiva.

Finalmente, como $f$ es inyectiva y sobreyectiva, entonces $f$ es biyectiva. Por lo que:
\begin{align*}
    \left|\left(\binom{[n+2]}{m}\right)\right| &= 
    \left(
        \binom{n+2}{m}
    \right)\\
    &= \left|
        \bigcup_{k=0}^m \left(\binom{[n+1]}{k}\right)
    \right|\\
    &= \sum_{k=0}^m \left|\left(\binom{[n+1]}{k}\right)\right|\\
    &= \sum_{k=0}^m \left(\binom{n+1}{k}\right)\\
\end{align*}
Concluyendo que:
\begin{align*}
    \sum_{k=0}^m\left(\binom{n+1}{k}\right)=\left(\binom{n+2}{m}\right)
\end{align*}
\textbf{Demostración 2.}

Una forma mas sencilla de demostrar la identidad es usando la regla de la adición, por medio de considerar la siguiente partición de las $k$-combinaciones del multiconjunto:
\begin{align*}
\left(\binom{n+2}{m}\right) = \{k_1\cdot 1, k_2\cdot 2, \ldots, k_{n+2}\cdot (n+2):k_i \in \mathbb N \text{ para todo $i$ } \text{ y } \sum_{i=1}^{n+2} k_i = m\}
\end{align*}
como $\{U_i\}_{k=1}^{m}$ con:
\begin{align*}
    U_k = \{\{k\cdot 1, k_2\cdot 2, \ldots, k_{n+2}\cdot (n+2)\}:k_i \in \mathbb N \text{ para todo $i$ } \text{ y } \sum_{i=2}^{n+2} k_i = m - k\}
\end{align*}
la cual cumple que:
\begin{itemize}
    \item 
    \begin{align*}
    \left|U_k\right| = \left(\binom{n+1}{k}\right) \quad k = 1, \ldots, n+2
\end{align*}
puesto que si fijamos el numero de $1$'s como $k$ entonces el cardinal de $U_k$ corresponderá a elegir $m-k$ combinaciones de los $n+1$ elementos restantes.
    \item $U_k \cap U_j = \emptyset$ para $k \neq j$. Dado que si asumimos sin perdida de generalidad $k>j$ entonces toda combinación de $U_k$ tiene al menos un $1$ mas que toda combinación de $U_j$.
    \item 
    \begin{align*}
        \bigcup_{k=0}^m U_k &= \left(\binom{[n+2]}{m}\right) \\
    \end{align*}
    Pues:
    \begin{itemize}
        \item [$\subseteq$] Si $X \in \bigcup_{k=0}^m U_k$ entonces $X \in U_k$ para algún $k$ y por lo tanto $X = \{k \cdot 1, k_2\cdot 2, \ldots, k_{n+2}\cdot (n+2)\}$ con $k+ \sum_{i=2}^{n+2} k_i = m$ y por lo tanto $X \in \left(\binom{[n+2]}{m}\right)$.
        \item [$\supseteq$] Si $X \in \left(\binom{[n+2]}{m}\right)$ entonces $X = \{k_1\cdot 1, k_2\cdot 2, \ldots, k_{n+2}\cdot (n+2)\}$ con $\sum_{i=2}^{n+2}  = m - k_1$ y por lo tanto $X \in U_{k_1}$ con $k_1\leq m$. Y por lo tanto $X \in \bigcup_{k=0}^m U_k$.
    \end{itemize}
\end{itemize} 

Asi por el principio de la adición:
\begin{align*}
     \left|\left(\binom{n+2}{m}\right) \right| &=
    \left|\left(\binom{[n+2]}{m}\right) \right|\\&= \left|\bigcup_{k=0}^m U_k\right|\\
    &= \sum_{k=0}^m \left|U_k\right|\\
    &= \sum_{k=0}^m \left(\binom{n+1}{k}\right)
\end{align*}

\end{solucion}